\documentclass{report}
\usepackage{hyperref}
% WARNING: THIS SHOULD BE MODIFIED DEPENDING ON THE LETTER/A4 SIZE
\oddsidemargin 0cm
\evensidemargin 0cm
\marginparsep 0cm
\marginparwidth 0cm
\parindent 0cm
\textwidth 16.5cm

\ifpdf
  \usepackage[pdftex]{graphicx}
\else
  \usepackage[dvips]{graphicx}
\fi

% definitons for warning and note tag
\usepackage[most]{tcolorbox}
\newtcolorbox{tcbwarning}{
 breakable,
 enhanced jigsaw,
 top=0pt,
 bottom=0pt,
 titlerule=0pt,
 bottomtitle=0pt,
 rightrule=0pt,
 toprule=0pt,
 bottomrule=0pt,
 colback=white,
 arc=0pt,
 outer arc=0pt,
 title style={white},
 fonttitle=\color{black}\bfseries,
 left=8pt,
 colframe=red,
 title={Warning:},
}
\newtcolorbox{tcbnote}{
 breakable,
 enhanced jigsaw,
 top=0pt,
 bottom=0pt,
 titlerule=0pt,
 bottomtitle=0pt,
 rightrule=0pt,
 toprule=0pt,
 bottomrule=0pt,
 colback=white,
 arc=0pt,
 outer arc=0pt,
 title style={white},
 fonttitle=\color{black}\bfseries,
 left=8pt,
 colframe=yellow,
 title={Note:},
}

\begin{document}
% special variable used for calculating some widths.
\newlength{\tmplength}
\chapter{Unit ok{\_}generic{\_}multiple{\_}constraints}
\section{Description}
Test that generics with types separated by semicolons are parsed correctly.\hfill\vspace*{1ex}

 See https://docwiki.embarcadero.com/RADStudio/Sydney/en/Constraints{\_}in{\_}Generics: """When you specify constraints, you separate multiple type parameters by semicolons, as you do with a parameter list declaration: """
\section{Overview}
\begin{description}
\item[\texttt{\begin{ttfamily}TMyClass\end{ttfamily} Class}]
\item[\texttt{\begin{ttfamily}TMyClassWithConstraints\end{ttfamily} Class}]
\item[\texttt{\begin{ttfamily}TMyClassWithConstraints2\end{ttfamily} Class}]
\end{description}
\section{Classes, Interfaces, Objects and Records}
\subsection*{TMyClass Class}
\subsubsection*{\large{\textbf{Hierarchy}}\normalsize\hspace{1ex}\hfill}
TMyClass {$>$} TObject
%%%%Description
\subsubsection*{\large{\textbf{Methods}}\normalsize\hspace{1ex}\hfill}
\paragraph*{Call}\hspace*{\fill}

\begin{list}{}{
\settowidth{\tmplength}{\textbf{Description}}
\setlength{\itemindent}{0cm}
\setlength{\listparindent}{0cm}
\setlength{\leftmargin}{\evensidemargin}
\addtolength{\leftmargin}{\tmplength}
\settowidth{\labelsep}{X}
\addtolength{\leftmargin}{\labelsep}
\setlength{\labelwidth}{\tmplength}
}
\begin{flushleft}
\item[\textbf{Declaration}\hfill]
\begin{ttfamily}
public function Call{$<$}A;R{$>$}(const Arg: A): R; overload;\end{ttfamily}


\end{flushleft}
\par
\item[\textbf{Description}]
Call with two type params separated by semicolon

\end{list}
\paragraph*{Call}\hspace*{\fill}

\begin{list}{}{
\settowidth{\tmplength}{\textbf{Description}}
\setlength{\itemindent}{0cm}
\setlength{\listparindent}{0cm}
\setlength{\leftmargin}{\evensidemargin}
\addtolength{\leftmargin}{\tmplength}
\settowidth{\labelsep}{X}
\addtolength{\leftmargin}{\labelsep}
\setlength{\labelwidth}{\tmplength}
}
\begin{flushleft}
\item[\textbf{Declaration}\hfill]
\begin{ttfamily}
public function Call{$<$}A;B;R{$>$}(const Arg1: A; const Arg2: B): R; overload;\end{ttfamily}


\end{flushleft}
\par
\item[\textbf{Description}]
Call with three type params separated by semicolons

\end{list}
\paragraph*{Execute}\hspace*{\fill}

\begin{list}{}{
\settowidth{\tmplength}{\textbf{Description}}
\setlength{\itemindent}{0cm}
\setlength{\listparindent}{0cm}
\setlength{\leftmargin}{\evensidemargin}
\addtolength{\leftmargin}{\tmplength}
\settowidth{\labelsep}{X}
\addtolength{\leftmargin}{\labelsep}
\setlength{\labelwidth}{\tmplength}
}
\begin{flushleft}
\item[\textbf{Declaration}\hfill]
\begin{ttfamily}
public procedure Execute{$<$}A;B{$>$}(const Arg1: A; const Arg2: B); overload;\end{ttfamily}


\end{flushleft}
\par
\item[\textbf{Description}]
Procedure variant

\end{list}
\subsection*{TMyClassWithConstraints Class}
\subsubsection*{\large{\textbf{Hierarchy}}\normalsize\hspace{1ex}\hfill}
TMyClassWithConstraints {$>$} \begin{ttfamily}TMyClass\end{ttfamily}(\ref{ok_generic_multiple_constraints.TMyClass}) {$>$} 
TObject
%%%%Description
\subsection*{TMyClassWithConstraints2 Class}
\subsubsection*{\large{\textbf{Hierarchy}}\normalsize\hspace{1ex}\hfill}
TMyClassWithConstraints2 {$>$} \begin{ttfamily}TMyClass\end{ttfamily}(\ref{ok_generic_multiple_constraints.TMyClass}) {$>$} 
TObject
%%%%Description
\section{Types}
\subsection*{TMethodFunc}
\begin{list}{}{
\settowidth{\tmplength}{\textbf{Description}}
\setlength{\itemindent}{0cm}
\setlength{\listparindent}{0cm}
\setlength{\leftmargin}{\evensidemargin}
\addtolength{\leftmargin}{\tmplength}
\settowidth{\labelsep}{X}
\addtolength{\leftmargin}{\labelsep}
\setlength{\labelwidth}{\tmplength}
}
\begin{flushleft}
\item[\textbf{Declaration}\hfill]
\begin{ttfamily}
TMethodFunc = function(Arg: A): R of object;\end{ttfamily}


\end{flushleft}
\par
\item[\textbf{Description}]
Generic type aliases with semicolons (in type section)

\end{list}
\subsection*{TMethodFunc2}
\begin{list}{}{
\settowidth{\tmplength}{\textbf{Description}}
\setlength{\itemindent}{0cm}
\setlength{\listparindent}{0cm}
\setlength{\leftmargin}{\evensidemargin}
\addtolength{\leftmargin}{\tmplength}
\settowidth{\labelsep}{X}
\addtolength{\leftmargin}{\labelsep}
\setlength{\labelwidth}{\tmplength}
}
\begin{flushleft}
\item[\textbf{Declaration}\hfill]
\begin{ttfamily}
TMethodFunc2 = function(Arg1: A; Arg2: B): R of object;\end{ttfamily}


\end{flushleft}
\end{list}
\end{document}
